\section{Completion of the repeated-exponent family}
\label{sec:completion}

The preceding factor-index reduction makes the repeated-exponent problem
finite at its first few levels. A uniform cubic sign interval handles every
later level, completing the family without a bound on either exponent.

\begin{theorem}
\label{thm:repeated}
For distinct primes $p,q,r$ and all integers $a,b\geq1$, the graph
$G_{p^a q^a r^b}$ is not Laplacian integral. The same conclusion holds
for every permutation of the exponent vector $(a,a,b)$.
\end{theorem}

The arithmetic reduction and a uniform sign obstruction for the cubic
combine to exhaust the family. We first settle one last finite factor level.

\begin{lemma}
\label{lem:third-index}
Every square-discriminant pair with index $j=3$ gives a nonintegral
Laplacian spectrum.
\end{lemma}
\begin{proof}
Proposition~\ref{prop:index} gives $d\mid K_3=320$ and $d=3a+4$.
For $a\geq2$, the complete divisor list of this shape is
$10,16,40,64,160$. The resulting values are
\begin{center}
\begin{tabular}{rrrrrl}
\toprule
$d$ & $a$ & $e$ & $v$ & $b$ & Outcome\\
\midrule
10 & 2 & 4 & 1 & 0 & $d>e$\\
16 & 4 & 36 & 7 & 0 & $b=0$\\
40 & 12 & 1080 & 83 & 4 & Root-free modulo $5$\\
64 & 20 & 5125 & 244 & 9 & Root-free modulo $7$\\
160 & 52 & 92274 & 1741 & 30 & Root-free modulo $7$\\
\bottomrule
\end{tabular}
\end{center}
The three positive-$b$ cubics are
\begin{align*}
f_{12,4}(x)&=x^3-568x^2+100032x-5248512,\\
f_{20,9}(x)&=x^3-1769x^2+992820x-174744000,\\
f_{52,30}(x)&=x^3-13394x^2+57771168x-80229951360.
\end{align*}
Their values at all residues of the indicated primes are, respectively,
$(3,3,3,4,2)$, $(3,2,4,1,6,4,1)$ and $(1,6,4,1,3,2,4)$.
Each monic cubic is therefore irreducible over $\mathbb Q$ and has
noninteger real roots. Section~\ref{sec:graph} supplies the lift to
noninteger graph eigenvalues.
\end{proof}

\begin{lemma}[A uniform consecutive-integer interval]
\label{lem:sign}
For $a,b\geq1$, put
\[
L(a)=\frac{a^3+2a^2+2a+1}{a(2a+1)}.
\]
If $b<L(a)$, then $f_{a,b}$ has a root strictly between $2ab-1$ and $2ab$.
\end{lemma}
\begin{proof}
Substitution in~\eqref{eq:cubic} gives
\begin{align*}
f_{a,b}(2ab)&=2a^2b^2>0,\\
f_{a,b}(2ab-1)&=-(a+b+1)
\bigl(a^3+2a^2+2a+1-a(2a+1)b\bigr).
\end{align*}
The second value is strictly negative when $b<L(a)$. The intermediate
value theorem gives the asserted noninteger, nonzero root.
\end{proof}

\begin{proof}[Proof of Theorem~\ref{thm:repeated}]
For $a=1$, Corollary~\ref{cor:quadratic} settles every positive $b$.
For $a\geq2$, nonsquare $D$ gives irrational antisymmetric eigenvalues.
When $D$ is square, index $j=0$ is Theorem~\ref{thm:boundary},
and indices $j=1,2$ are settled in Section~\ref{sec:bounds}.
Lemma~\ref{lem:third-index} settles $j=3$.

It remains to consider $j\geq4$. Then $d\geq Q=4a+5$, and
$e\geq d\geq Q$ implies
\[
\frac MQ+Q-(d+e)=(d-Q)\left(\frac eQ-1\right)\geq0.
\]
Since $d+e=(a+1)^2b+a(3a+1)$, we obtain
\[
b\leq B_4(a)=\frac{M/Q+Q-a(3a+1)}{(a+1)^2}
=\frac{(a-5)(2a-5)}{4a+5}.
\]
For every $a\geq2$,
\[
L(a)-B_4(a)=
\frac{41a^3-17a^2-11a+5}{a(2a+1)(4a+5)}>0.
\]
Indeed, with $z=a-2\geq0$, the numerator becomes
$41z^3+229z^2+413z+243$. Thus $b\leq B_4(a)<L(a)$, and
Lemma~\ref{lem:sign} gives a noninteger cubic root, which lifts across
the universal vertices. These cases exhaust all positive $a,b$.
Permuting the prime factors preserves the graph's isomorphism type.
\end{proof}

The preceding cutoffs and fixed-exponent theorems are intermediate
consequences of the arithmetic reduction. The uniform interval now
settles the infinite range left by those cutoffs, without a parameter scan.
